\section{Periodic finite parts and harmonic contact terms}
\label{per:sec:main}

We begin by consolidating the untwisted foundation already present in
three archived ProveIt continuations
\cite{RepoConvolutionAlgebra,RepoConvolutionCalculus,RepoShiftedJets}.
The endpoint normalization is fixed by analytic continuation from the
integrable Hurwitz family. Its spectral finite parts and argument
derivatives differ by an explicit harmonic-number multiple of a Dirac
mass derivative. This distinction will persist in the twisted extension.

The normalized Hurwitz convolution is classical
\cite[Remark 3.2]{per:Sun}; complex-power distributions provide related
background \cite[Section 5]{per:GaySebbar}. Scalar meromorphic
continuation after pairing with a smooth function is also covered by
Kim's Hurwitz-transform theorem
\cite[Theorem 1 and Corollary 2]{Kim2026}. The full proofs below establish
the conventions used in the rest of this article and independently check
the archived formulas; the untwisted Bell algebra and primitive tower are
not new-to-repository claims.

\subsection{Conventions and the meromorphic periodic family}

Write $\mathbb T=\mathbb R/\mathbb Z$, use normalized Lebesgue measure,
and put
\[
 \widehat U(n)=\langle U,e^{-2\pi i n x}\rangle,\qquad
 D=\frac{d}{dx},\qquad \Delta=\delta_0-1.
\]
The Dirac mass has Fourier coefficient $1$ at every integer. Thus
$\widehat\Delta(0)=0$ and $\widehat\Delta(n)=1$ for $n\ne0$.
On the convolution algebra of distributions with zero mean, $\Delta$
is the identity. Reflection is $\check U(x)=U(-x)$, so
$\widehat{\check U}(n)=\widehat U(-n)$. Throughout,
\begin{equation}\label{per:eq:L}
 L_n=\log(2\pi |n|)+\frac{\pi i}{2}\operatorname{sgn}n,
 \qquad n\ne0.
\end{equation}
Consequently $(2\pi i n)^z$ always means $\exp(zL_n)$.

\begin{theorem}[The entire normalized family and all residues]
\label{per:thm:family}
There is an entire distribution-valued function $W_s$ characterized by
\begin{equation}\label{per:eq:W}
 \widehat W_s(0)=0,\qquad
 \widehat W_s(n)=(2\pi i n)^{s-1}\quad(n\ne0).
\end{equation}
The family
\begin{equation}\label{per:eq:Z}
 Z_s=\Gamma(1-s)W_s
\end{equation}
is the unique meromorphic continuation, as a periodic distribution, of
$x\mapsto\zeta(s,x)$ from $\operatorname{Re}s<1$. Its poles are simple
and occur at every positive integer. Their residues are
\begin{align}
 \operatorname*{Res}_{s=1} Z_s&=1-\delta_0=-\Delta,
 \label{per:eq:res1}\\
 \operatorname*{Res}_{s=k+1} Z_s
 &=\frac{(-1)^{k+1}}{k!}\delta_0^{(k)}\qquad(k\ge1).
 \label{per:eq:resk}
\end{align}
For all complex $s,t$, the identities
\begin{equation}\label{per:eq:group}
 DW_s=W_{s+1},\qquad W_s*W_t=W_{s+t-1}
\end{equation}
hold. Equivalently, as meromorphic distribution identities,
\begin{equation}\label{per:eq:meroconv}
 DZ_s=-sZ_{s+1},\qquad
 Z_s*Z_t=\frac{\Gamma(1-s)\Gamma(1-t)}{\Gamma(2-s-t)}Z_{s+t-1}.
\end{equation}
At intersections of polar divisors the second expression is interpreted
by its equal Gamma-normalized expression, before coefficients are taken.
\end{theorem}

\begin{proof}
On every compact set of $s$ the coefficients in \eqref{per:eq:W} have
uniform polynomial growth in $|n|$. Every fixed order derivative in $s$
adds only a power of $L_n$. Since the Fourier coefficients of a smooth
test function decrease faster than every inverse power, the series
\[
 \langle W_s,\phi\rangle
 =\sum_{n\ne0}(2\pi i n)^{s-1}\widehat\phi(-n)
\]
and all its fixed order derivatives converge normally on compact sets.
This constructs the entire family. Hurwitz's classical Fourier formula
\cite[\S25.11]{DLMFHurwitz} identifies \eqref{per:eq:Z} with $\zeta(s,x)$
for $\operatorname{Re}s<0$. The relation
$\zeta(s,x)=x^{-s}+\zeta(s,1+x)$ shows that the latter is holomorphic
with values in $L^1(0,1)$ on $\operatorname{Re}s<1$, so the
identification extends throughout that half-plane. The identity theorem
for each test-function pairing gives uniqueness of the continuation.

At $s=k+1+w$,
\begin{equation}\label{per:eq:gammapole}
 \Gamma(-k-w)=\frac{(-1)^{k+1}}{k!w}
       +\frac{(-1)^k}{k!}(H_k-\gamma)+O(w),
 \qquad H_0=0.
\end{equation}
The Fourier coefficients show that $W_1=\Delta$ and
$W_{k+1}=D^k\Delta=\delta_0^{(k)}$ for $k\ge1$.
This proves the residues and also shows that none of the listed poles
is removable. Multiplication of Fourier coefficients proves
\eqref{per:eq:group}. Convolution of two distributions on the compact
group $\mathbb T$ is well defined. Normal convergence after pairing with
a test function justifies the parameter identities jointly. Finally,
$\Gamma(1-s)=-s\Gamma(-s)$ proves the derivative identity in
\eqref{per:eq:meroconv}.
\end{proof}

The residue at $s=1$ has both a constant and a Dirac term. The constant
is the ordinary Hurwitz pole seen on every open subinterval of $(0,1)$;
the Dirac term comes from the nonintegrable endpoint $x^{-s}$. Omitting
either term destroys the zero Fourier coefficient. The poles at
$2,3,\ldots$ are entirely supported at the endpoint and are invisible
in the pointwise meromorphic function of $s$ for a fixed $0<x<1$.

\subsection{Cutoff definitions and the normalization at the endpoint}

Define $T_n$ by the Laurent expansion
\begin{equation}\label{per:eq:Tdef}
 Z_{1+w}=-\frac{\Delta}{w}
       +\sum_{n=0}^{\infty}\frac{(-1)^n}{n!}T_nw^n.
\end{equation}
Each $T_n$ agrees with $\gamma_n(x)$ on $0<x<1$. Define also
\begin{equation}\label{per:eq:Pdef}
 F_k=\operatorname{FP}_{s=k+1}Z_s,\qquad
 P_k=(-1)^{k+1}k!F_k\quad(k\ge0),
\end{equation}
where $\operatorname{FP}$ denotes the constant Laurent coefficient.
Thus $F_0=T_0$ and $P_0=-T_0$. For $k\ge1$, the usual
$\psi^{(k)}(x)=(-1)^{k+1}k!\zeta(k+1,x)$ shows that $P_k$ agrees with
the ordinary $k$th polygamma function away from $0$.

\begin{proposition}[Equivalent Hadamard finite parts]
\label{per:prop:cutoffs}
For every smooth periodic test function $\phi$,
\begin{equation}\label{per:eq:Tcutoff}
 \langle T_n,\phi\rangle
 =\lim_{\epsilon\downarrow0}\left\{
   \int_\epsilon^1\gamma_n(x)\phi(x)\,dx
    +\frac{\phi(0)}{n+1}\log^{n+1}\epsilon\right\}.
\end{equation}
For $k\ge1$,
\begin{align}
 \langle F_k,\phi\rangle
 =\lim_{\epsilon\downarrow0}\biggl\{&
   \int_\epsilon^1\zeta(k+1,x)\phi(x)\,dx
   -\sum_{j=0}^{k-1}\frac{\phi^{(j)}(0)}{j!}
                       \frac{\epsilon^{j-k}}{k-j}
   +\frac{\phi^{(k)}(0)}{k!}\log\epsilon\biggr\}.
 \label{per:eq:Fcutoff}
\end{align}
All these distributions have zero mean.
\end{proposition}

\begin{proof}
Subtract a Taylor polynomial of the test function at $0$. For any
$N\ge0$, the pairing of $Z_s$ has the continuation
\begin{align}
 \langle Z_s,\phi\rangle={}&
 \int_0^1x^{-s}\left(\phi(x)-\sum_{j=0}^{N}
                         \frac{\phi^{(j)}(0)}{j!}x^j\right)dx
 +\sum_{j=0}^{N}\frac{\phi^{(j)}(0)}{j!(j+1-s)}\notag\\
 &+\int_0^1\zeta(s,1+x)\phi(x)\,dx,
 \label{per:eq:localcontinuation}
\end{align}
in $\operatorname{Re}s<N+2$, apart from the displayed and Hurwitz
poles. Taking the constant coefficient at $s=k+1$, with $N=k$,
proves \eqref{per:eq:Fcutoff} by elementary integration of the Taylor
monomials. At $s=1+w$, use
\[
 \gamma_n(x)=\frac{\log^n x}{x}+\gamma_n(1+x).
\]
The coefficient in \eqref{per:eq:localcontinuation} is
\[
 \int_0^1\frac{\log^n x}{x}\bigl(\phi(x)-\phi(0)\bigr)dx
 +\int_0^1\gamma_n(1+x)\phi(x)\,dx.
\]
This is precisely \eqref{per:eq:Tcutoff}. The zero-mean assertion
follows from \eqref{per:eq:W} by taking Laurent coefficients.
\end{proof}

Zero mean alone does not specify a periodic extension of a singular
function: derivatives of a Dirac mass have zero mean as well. Here the
normalization is fixed by \eqref{per:eq:Z}, or equivalently by the
cutoff formulas. This qualification is essential in using the word
``canonical.''

\subsection{The harmonic contact correction}

\begin{theorem}[Spectral finite parts and compatible derivatives]
\label{per:thm:contact}
For $n\ne0$ and $k\ge0$,
\begin{equation}\label{per:eq:Pfourier}
 \widehat P_k(n)=(2\pi i n)^k\bigl(\gamma+L_n-H_k\bigr),
 \qquad \widehat P_k(0)=0.
\end{equation}
Consequently,
\begin{equation}\label{per:eq:anomaly}
 DP_k=P_{k+1}+\frac{1}{k+1}\delta_0^{(k+1)}.
\end{equation}
If
\[
 g(x)=\log\Gamma(x)-\frac12\log(2\pi),\quad 0<x<1,
\]
is periodized as a locally integrable function, then $Dg=P_0$. The
derivative-compatible extensions
\begin{equation}\label{per:eq:Q}
 Q_k=D^{k+1}g=D^kP_0
\end{equation}
satisfy $DQ_k=Q_{k+1}$ exactly, and
\begin{equation}\label{per:eq:Qcorrection}
 Q_k=P_k+H_k\delta_0^{(k)}\qquad(k\ge1),\qquad Q_0=P_0.
\end{equation}
\end{theorem}

\begin{proof}
Expand $W_{k+1+w}=W_{k+1}+w\partial_sW_s|_{s=k+1}+O(w^2)$
against \eqref{per:eq:gammapole}; after multiplication by
$(-1)^{k+1}k!$, the finite Fourier coefficient is
$(2\pi i n)^k(\gamma+L_n-H_k)$.
Since $H_{k+1}-H_k=1/(k+1)$, differentiation of this coefficient gives
\eqref{per:eq:anomaly}. Lerch's identity gives $g=\partial_sZ_s|_{s=0}$.
Taking the first coefficient in $DZ_s=-sZ_{s+1}$ at $s=0$ gives
$Dg=-T_0=P_0$; the coefficient at degree zero separately says
$DZ_0=\Delta$. Finally the Fourier coefficient of $D^kP_0-P_k$ is
$H_k(2\pi i n)^k$, which is \eqref{per:eq:Qcorrection}.
\end{proof}

For example, the extension given by the finite part of
$\zeta(s,x)$ at $s=2$ yields a periodic trigamma $P_1$, whereas
$D^2g=P_1+\delta_0'$. Both agree with $\psi'(x)$ on $(0,1)$.
The discrepancy is supported precisely at the identified endpoints;
it is not a failure of ordinary differentiation on the open interval.

